% 从提问到播报的前后端状态机
\begin{tikzpicture}[x=1mm,y=1mm]
  \path[use as bounding box] (0,0) rectangle (158,63);
  \tikzset{
    hhState/.style={
      draw=JinmenBlue, fill=PaleJade!52, line width=.72pt,
      rounded corners=1.2mm, minimum width=22mm, minimum height=14mm,
      text width=18.5mm, align=center, inner sep=1.1mm,
      font=\small\sffamily, text=Ink
    },
    hhStateGuard/.style={hhState,draw=Cinnabar,fill=RicePaper},
    hhStateFlow/.style={
      -{Latex[length=2mm,width=1.35mm]}, draw=JinmenBlue,
      line width=.72pt, line cap=round, line join=round
    },
    hhStateFail/.style={
      -{Latex[length=2mm,width=1.35mm]}, draw=Cinnabar,
      line width=.72pt, line cap=round, line join=round
    }
  }

  \node[hhState] (idle) at (11,48) {\textbf{待机}\\[-.2ex]\footnotesize 接收问题};
  \node[hhState] (retrieve) at (38,48) {\textbf{检索}\\[-.2ex]\footnotesize 候选证据};
  \node[hhStateGuard] (gate) at (65,48) {\textbf{证据门控}\\[-.2ex]\footnotesize 范围・相关度};
  \node[hhState] (generate) at (92,48) {\textbf{生成}\\[-.2ex]\footnotesize 事实段};
  \node[hhState] (verify) at (119,48) {\textbf{引用复核}\\[-.2ex]\footnotesize ID・数字・语义};
  \node[hhState] (speech) at (146,48) {\textbf{播报}\\[-.2ex]\footnotesize 一次性令牌};

  \node[hhStateGuard,minimum width=30mm,text width=26mm] (scopeRefuse) at (65,9)
    {\textbf{范围拒答}\\[-.2ex]\scriptsize 越界或证据不足};
  \node[hhStateGuard,minimum width=30mm,text width=26mm] (checkRefuse) at (119,9)
    {\textbf{校验拒答}\\[-.2ex]\scriptsize 引用或语义不通过};
  \node[hhStateGuard,minimum width=22mm,text width=18mm] (error) at (146,9)
    {\textbf{可见错误}\\[-.2ex]\scriptsize 不伪造结果};

  \draw[hhStateFlow] (idle.east) -- (retrieve.west);
  \draw[hhStateFlow] (retrieve.east) -- (gate.west);
  \draw[hhStateFlow] (gate.east) -- node[above,font=\tiny\sffamily,text=BrickGrey] {是} (generate.west);
  \draw[hhStateFlow] (generate.east) -- (verify.west);
  \draw[hhStateFlow] (verify.east) -- node[above,font=\tiny\sffamily,text=BrickGrey] {通过} (speech.west);

  % 两条拒答分支均垂直进入各自终止状态，不交叉、不共用轨道。
  \draw[hhStateFail] (gate.south) -- node[right,font=\tiny\sffamily,text=Cinnabar] {否} (scopeRefuse.north);
  \draw[hhStateFail] (verify.south) -- node[right,font=\tiny\sffamily,text=Cinnabar] {否} (checkRefuse.north);

  % 云端异常是任一云调用均可能触发的独立终止条件，不错误归属于某一步。
  \node[font=\tiny\sffamily,text=Cinnabar,align=center] (cloudFail) at (146,27)
    {任一云调用失败};
  \draw[hhStateFail] (cloudFail.south) -- (error.north);

  \node[font=\scriptsize\sffamily,text=BrickGrey,anchor=south]
    at (80,59) {SSE 仅用于传输当前服务端状态；新问题会取消旧请求};
\end{tikzpicture}
